% ============================================================
% conclusion.tex - Заключение и литература
% ============================================================
% !TeX program = xelatex

\section*{Заключение}
\addcontentsline{toc}{section}{Заключение}

В данном учебно-методическом пособии мы прошли полный путь от «сырых» физических данных ТГц-спектроскопии до прецизионного восстановления геометрических параметров микроструктур. Реализованный нами комплекс модулей и интерактивное графическое приложение связывают фундаментальную электродинамику с современной практикой цифровой обработки сигналов на языке Python.

\subsection*{Основные итоги работы}
В ходе выполнения лабораторной работы и освоения материала пособия были решены следующие задачи:
\begin{itemize}
    \item \textbf{Загрузка и организация данных:} Разработан гибкий загрузчик экспериментальных файлов, реализующий СУБД-подобное хранилище в оперативной памяти для исключения хаоса при работе с множеством угловых и повторных измерений.
    \item \textbf{Цифровая обработка сигналов:} Реализован конвейер DSP, включающий удаление постоянного смещения поля, Гауссово временное окно для фильтрации паразитных отражений волновода и БПФ для перехода в частотную область.
    \item \textbf{Электродинамическое моделирование:} Реализован расчет энергетического пропускания системы поляризаторов по строгой модели Бланко и матричному формализму Джонса для произвольных углов поворота ротатора.
    \item \textbf{Численная оптимизация:} Автоматизирована процедура локальной и пакетной спектральной подгонки параметров решетки методом Нелдера-Мида, позволившая доказать адекватность модели по постоянству геометрических параметров ($P$ и $D$) в широком диапазоне частот.
\end{itemize}

\subsection*{Полученные навыки}
В результате изучения пособия студент приобретает ключевые компетенции на стыке физики и программирования:
\begin{itemize}
    \item Навыки проектирования структурированных, поддерживаемых научных библиотек на языке Python с использованием NumPy и SciPy.
    \item Понимание физического смысла спектральных преобразований, фильтрации и принципа неопределенности Фурье на реальных экспериментальных сигналах.
    \item Опыт создания многооконных интерактивных GUI-интерфейсов с обратной связью на базе виджетов Matplotlib без привлечения тяжеловесных фреймворков.
    \item Практические навыки применения численных алгоритмов оптимизации к решению обратных физических задач и интерпретации их результатов.
\end{itemize}

\subsection*{Направления дальнейшего развития}
Созданный комплекс представляет собой открытую систему, готовую к расширению. В качестве тем для курсовых или исследовательских работ предлагаются следующие направления:
\begin{itemize}
    \item \textbf{Учет конечной проводимости проволок:} Переход от модели идеального проводника Бланко к модели Манабэ, учитывающей скин-эффект в вольфрамовой или золотой проволоке на частотах выше $1.5$ ТГц.
    \item \textbf{Моделирование нерегулярности намотки:} Учет случайных микро-отклонений шага решетки, приводящих к диффузному рассеянию коротких волн.
    \item \textbf{Экспорт результатов:} Разработка модуля автоматического сохранения результатов оптимизации в форматы CSV или JSON для накопления базы данных исследуемых образцов.
\end{itemize}

\subsection*{Литература}
\begin{thebibliography}{9}
	\bibitem{blanco1986} Blanco A., Fonti S., Piacente A. Transmission coefficients of free-standing wire grids in the far infrared: a theoretical approach for easy computation. – Infrared Physics, 1986, vol. 26, no. 6, pp. 357-362.
	\bibitem{manabe2005} Manabe T., Murk A. Transmission and reflection characteristics of slightly irregular wire-grids with finite conductivity for arbitrary angles of incidence and grid rotation. – IEEE Transactions on Antennas and Propagation, 2005, vol. 53, no. 1, pp. 250-259.
	\bibitem{chambers1980} Chambers W.G., Mok C.L., Parker T.J. Theory of the scattering of electromagnetic waves by a regular grid of parallel cylindrical wires with circular cross section. – Journal of Physics A: Mathematical and General, 1980, vol. 13, pp. 1433-1441.
\end{thebibliography}
